\section{Conclusion}
\label{sec:conclusion}

This work presents a systematic evaluation of symbolic constraint aggregation for LLM majority voting across five models, four reasoning domains, and more than 1,300 instances. The central finding is negative: when an LLM extracts first-order logic constraints from its own reasoning traces and a MAX-SAT solver aggregates them, accuracy does not improve over standard majority voting. The gap ranges from $-3.3$ to $+2.0$ percentage points, with no condition reaching statistical significance. We trace this failure to a capability trap: self-extracted constraints are semantically coupled to the answers they accompany and carry no independent discriminative signal. The MAX-SAT solver agrees with simple constraint counting on 803 of 804 instances, which confirms that satisfiability structure adds nothing when constraints are redundant with votes.

An oracle control experiment isolates constraint quality as the sole bottleneck. Replacing self-extracted constraints with gold FOL annotations yields $+8.82$ percentage points over majority voting ($p < 0.001$), while all other pipeline components remain fixed. The aggregation mechanism is sound; the problem lies entirely in same-model self-extraction. This result quantifies an $8.82$-point gap that remains recoverable through higher-quality constraint sources.

For practitioners, we identify the multi-candidate rate $C_0$ as a zero-cost diagnostic: when $C_0$ falls below 5\%, constraint-based methods cannot improve over majority voting regardless of constraint quality, and investment in extraction is unwarranted. For the field, our results suggest that productive directions include trained FOL extractors, external theorem provers, and cross-model constraint extraction, all of which break the circular dependency between generation and verification that defines the capability trap.
